%----------------------------------------------------------------------------------------
% Preambulo y Configuración
%----------------------------------------------------------------------------------------

\documentclass[
    11pt,
    spanish,
    singlespacing,
    parskip,
    headsepline,
    bookmarks=true,
    unicode=true,
    pdftoolbar=true,
    pdfmenubar=true,
    pdffitwindow=false,
    colorlinks=true,
    linkcolor=blue,
    citecolor=blue,
    urlcolor=blue
]{MastersDoctoralThesis}

\usepackage[utf8]{inputenc} % Codificación de entrada UTF-8
\usepackage[T1]{fontenc}    % Codificación de salida para caracteres especiales
\usepackage{graphicx}       % Manejo de gráficos
\usepackage{eso-pic}        % Permite agregar fondos
\usepackage{hyperref}       % Manejo de hipervínculos y marcadores

% Redefinición de caracteres problemáticos en marcadores
\hypersetup{
    pdftitle={Título del Documento},
    pdfauthor={Autor del Documento},
    pdfkeywords={Sistemas Embebidos, Internet de las Cosas, Inteligencia Artificial},
    pdfstartview={FitH},
    unicode=true,
    colorlinks=true,
    linkcolor=blue,
    citecolor=blue,
    urlcolor=blue
}

\pdfstringdefDisableCommands{%
  \def\texttt#1{#1}%
  \def\textbf#1{#1}%
  \def\textit#1{#1}%
  \def\"{\"}%
  \def\~{~}%
  \def\'{'}%
  \def\^{}%
  \def\textunderscore{\_} % Manejo del subrayado en marcadores
}


% Definir comandos requeridos por la clase
\newcommand{\degreename}{Maestría en Ciencias} % Cambia según tu título
\newcommand{\univname}{Universidad Nacional de Ejemplo} % Cambia según tu universidad
\newcommand{\keywordnames}{Palabras clave:}
%----------------------------------------------------------------------------------------
% Documento Principal
%----------------------------------------------------------------------------------------

\begin{document}

% Configuración de la portada
\posgrado{Carrera / Maestría}
\keywords{Sistemas Embebidos, Internet de las Cosas, Inteligencia Artificial}

% Incluir la portada desde un archivo separado
\include{portada}

% Configuración del contenido preliminar
\frontmatter % Usar numeración romana para las páginas preliminares
\pagestyle{plain} % Estilo de encabezado simple

%----------------------------------------------------------------------------------------
% Resumen
%----------------------------------------------------------------------------------------

\begin{abstract}
\addchaptertocentry{\abstractname} % Agregar resumen al índice
En este trabajo se desarrolló un modelo de aprendizaje automático, junto con un flujo de procesamiento de datos reproducible, para proyectar las contrataciones de una empresa a partir de fotografías mensuales y trimestrales históricas de su sistema de gestión comercial. El propósito fue reemplazar las estimaciones manuales y subjetivas utilizadas hasta el momento por proyecciones basadas en evidencia, reproducibles y trazables, que sirvieran de apoyo a la planificación financiera de la organización. Para ello se construyó un conjunto de datos respetando la correcta separación temporal de la información, de modo de evitar que el modelo utilizara datos posteriores a la fecha de cada proyección, se compararon distintos enfoques de modelado, que abarcaron regresión, series temporales y redes neuronales, bajo una metodología de validación temporal común, y se seleccionó y ajustó el de mejor desempeño. El modelo resultante redujo el error de proyección y superó de manera consistente al método manual tomado como referencia.
La memoria presenta el contexto y la formulación del problema, el marco teórico, el diseño e implementación del flujo de trabajo y de los modelos, los ensayos realizados con sus resultados, y las conclusiones. En el proyecto se aplicaron de manera integrada contenidos de la especialización, entre ellos el análisis y la preparación de datos, el aprendizaje automático, las series temporales y la evaluación de modelos.
\end{abstract}

%----------------------------------------------------------------------------------------
% Agradecimientos
%----------------------------------------------------------------------------------------

%\begin{acknowledgements}
%\vspace{1.5cm}
%Esta sección es para agradecimientos personales y es totalmente \textbf{OPCIONAL}.
%\end{acknowledgements}

%----------------------------------------------------------------------------------------
% Índice
%----------------------------------------------------------------------------------------

\tableofcontents
\listoffigures
\listoftables

%----------------------------------------------------------------------------------------
% Dedicatoria
%----------------------------------------------------------------------------------------

%\dedicatory{\textbf{Dedicado a... [OPCIONAL]}}

%----------------------------------------------------------------------------------------
% Capítulos
%----------------------------------------------------------------------------------------

\mainmatter % Iniciar numeración numérica para el contenido principal
\pagestyle{thesis} % Estilo de encabezado de tesis

% Incluir capítulos desde archivos separados
% Chapter 1

\chapter{Introducción general}

\label{Chapter1}
\label{IntroGeneral}

En este capítulo se introduce el problema de la proyección de contrataciones a partir de información comercial histórica. Se presenta el rol del pronóstico de ventas en la planificación de una organización y se describen los desafíos propios de las ventas entre empresas. Luego se expone el contexto que motivó el trabajo y se formula el problema técnico. Finalmente, se revisan los principales enfoques propuestos en la literatura y se establecen los objetivos y el alcance del trabajo.

\section{Desafíos del pronóstico de ventas B2B}
\label{sec:desafiosPronosticoB2B}

La proyección de ventas constituye un insumo central para la planificación financiera de una organización. A partir de ella se estiman ingresos, se asignan recursos, se definen metas y se anticipan necesidades de capacidad. En entornos empresariales, esta tarea suele combinar información histórica con el conocimiento de los responsables comerciales. Los sistemas de pronóstico brindan una base cuantitativa, mientras que las personas incorporan información contextual que no siempre se encuentra representada en los datos. Sin embargo, los ajustes basados en juicio experto no mejoran necesariamente la precisión y pueden introducir sesgos, en particular cuando son frecuentes, pequeños o excesivamente optimistas \citep{Fildes2009}.

Los sistemas de gestión de relaciones con clientes o CRM (\textit{Customer Relationship Management}) registran de manera estructurada la actividad comercial. Para cada oportunidad pueden almacenar su monto esperado, la probabilidad estimada de cierre, la etapa dentro del proceso de venta, la fecha prevista de concreción y otras características relevantes. Además de facilitar el seguimiento operativo, estos registros permiten reconstruir cómo evolucionó el \textit{pipeline} comercial y ofrecen una fuente de datos para desarrollar pronósticos reproducibles.

En mercados de empresa a empresa o B2B (\textit{business to business}), el modelado presenta dificultades particulares. La cantidad de operaciones suele ser menor que en contextos de consumo masivo, los montos son heterogéneos y unas pocas oportunidades de gran tamaño pueden modificar de manera significativa el total de un período. También son frecuentes los cambios en las fechas esperadas de cierre y en las valoraciones subjetivas del equipo comercial. Se han propuesto modelos de aprendizaje automático para estimar la probabilidad de éxito de oportunidades B2B a partir de la evolución del \textit{pipeline} \citep{Yan2015}. Otros trabajos utilizaron información comercial anticipada para mejorar procesos manuales de pronóstico en casos reales de ventas B2B \citep{Rohaan2022}.

El trabajo aborda una variante diferente y complementaria. En lugar de estimar únicamente la probabilidad de éxito de cada oportunidad, busca proyectar el volumen agregado de contrataciones, entendido como el monto de ventas efectivamente comprometidas durante un horizonte futuro. La salida sirve como apoyo para la planificación financiera trimestral y permite una comparación directa con el pronóstico utilizado por la organización.

\section{Contexto y motivación}
\label{sec:contextoMotivacion}

OceanSound Partners es una firma de inversión en empresas de tecnología y servicios habilitados por tecnología que atienden mercados gubernamentales y empresariales \citep{OceanSoundWebsite}. Una de las empresas de su portafolio administra la actividad comercial mediante un CRM en el que registra la evolución de sus oportunidades de venta. La planificación financiera depende, en buena medida, de la proyección de las contrataciones futuras.

Antes de este trabajo, esas proyecciones se elaboraban a partir del pronóstico almacenado en el CRM y de la evaluación manual del equipo comercial. El proceso dependía del criterio de las personas involucradas, era difícil de reproducir y no contaba con un procedimiento común para medir su error histórico. Esta situación dificultaba anticipar desvíos en los ingresos, asignar recursos con suficiente previsión y distinguir si una corrección provenía de nueva información comercial o de una apreciación subjetiva. Como consecuencia, la organización tenía menos capacidad para detectar de manera temprana necesidades de contratación, ajustes presupuestarios y oportunidades que requerían atención.

El CRM conserva una fotografía semanal, denominada \textit{snapshot}, del estado completo del \textit{pipeline}. Una misma oportunidad aparece en sucesivos \textit{snapshots} con la información disponible en cada fecha. Esta estructura permite estudiar no solo el resultado final de una venta, sino también la evolución de las variables que podían observarse antes de que ese resultado ocurriera.

La motivación del trabajo surge de dos oportunidades. Desde la perspectiva organizacional, se busca disponer de un pronóstico cuantitativo, trazable y comparable con el método vigente. Desde la perspectiva académica, el problema permite integrar preparación y análisis de datos, aprendizaje automático, series temporales, validación de modelos y gestión de experimentos sobre un caso real.

\section{Descripción del problema}
\label{sec:descripcionProblema}

El conjunto de datos tiene estructura de panel longitudinal: cada fila representa el estado de una oportunidad en una fecha de observación y una misma oportunidad se repite a lo largo de múltiples \textit{snapshots}. Las contrataciones se determinan a partir de las oportunidades ganadas y de su fecha efectiva de cierre. Por esta razón, una agregación directa de todas las filas produciría duplicaciones y mezclaría estados conocidos en momentos diferentes.

El desafío metodológico principal consiste en respetar la disponibilidad temporal de la información. Para producir una proyección desde un instante $T$, las variables predictoras solo pueden utilizar datos que ya eran conocidos en ese instante. Incorporar el estado final de una oportunidad, su fecha de cierre efectiva o una transformación calculada con observaciones posteriores constituiría una fuga de información o \textit{data leakage}. Esta clase de error genera métricas artificialmente optimistas y un modelo que no puede reproducir su desempeño cuando se utiliza sobre períodos futuros.

La fuente combina así una gran cantidad de filas con una historia temporal relativamente corta. El problema requiere transformar los estados históricos del \textit{pipeline} en observaciones comparables, definir un monto de contrataciones sin contar una misma oportunidad más de una vez y relacionar cada observación con un período futuro. También exige considerar que los meses recientes pueden contener oportunidades cuyo resultado definitivo todavía no se encontraba registrado. El diseño adoptado para resolver estas dificultades se presenta en el capítulo~\ref{Chapter3}.

\section{Estado del arte}
\label{sec:estadoArte}

El pronóstico de ventas puede abordarse mediante métodos estadísticos, modelos de aprendizaje automático, redes neuronales y estimaciones basadas en juicio experto. No existe un enfoque universalmente superior: el resultado depende de la cantidad de datos, su frecuencia, el horizonte de decisión, la estabilidad del proceso y la disponibilidad de variables explicativas.

\subsection{Pronósticos estadísticos y basados en juicio experto}

Los modelos autorregresivos integrados de media móvil o ARIMA (\textit{Autoregressive Integrated Moving Average}) constituyen una referencia clásica para representar autocorrelación, tendencia y estacionalidad \citep{BoxJenkins2015}. Su formulación es interpretable y puede resultar efectiva cuando la estructura histórica de la serie permanece relativamente estable. Como limitación para este trabajo, un modelo puramente univariado no utiliza directamente el estado actual de las oportunidades comerciales.

Prophet propone un modelo descomponible con tendencia, estacionalidades y eventos, acompañado por parámetros que pueden ser ajustados por un analista \citep{TaylorLetham2018}. Fue concebido para facilitar el pronóstico de numerosas series de negocio y ofrece una referencia intermedia entre automatización e intervención humana. No obstante, tanto Prophet como ARIMA necesitan suficientes observaciones para identificar patrones temporales de manera confiable.

Los pronósticos basados en criterio comercial incorporan conocimiento que puede no estar codificado en los datos, pero su desempeño depende de cómo se realicen los ajustes. La evidencia muestra que algunos ajustes mejoran la precisión, mientras que otros la reducen y presentan sesgo optimista \citep{Fildes2009}. Por ese motivo, resulta necesario medir el valor agregado por la intervención humana y conservar el pronóstico del CRM como una referencia explícita.

\subsection{Aprendizaje automático aplicado al embudo comercial}

Los métodos basados en árboles permiten capturar relaciones no lineales e interacciones sin imponer una forma funcional única. Los bosques aleatorios o \textit{Random Forest} reducen la varianza mediante la combinación de árboles construidos sobre muestras y subconjuntos aleatorios de variables \citep{Breiman2001}. XGBoost agrega árboles de manera secuencial mediante \textit{gradient boosting} e incorpora regularización y optimizaciones para datos dispersos \citep{ChenGuestrin2016}. Ambos enfoques pueden utilizar rezagos históricos junto con medidas agregadas del \textit{pipeline}, aunque requieren control de complejidad cuando la cantidad de períodos es pequeña.

Las redes recurrentes LSTM (\textit{Long Short Term Memory}) fueron diseñadas para aprender dependencias prolongadas mediante mecanismos de memoria y compuertas \citep{HochreiterSchmidhuber1997}. Su flexibilidad permite representar dinámicas temporales complejas, pero implica una cantidad mayor de parámetros y una necesidad de datos que puede resultar desfavorable en series empresariales cortas. En este contexto, su inclusión constituye una hipótesis que debe evaluarse y no una garantía de mejora.

En el dominio B2B se utilizaron modelos para estimar la probabilidad de éxito de oportunidades a partir de información temporal del \textit{pipeline} \citep{Yan2015}. También se aplicó aprendizaje supervisado a información comercial anticipada para complementar procesos manuales de pronóstico \citep{Rohaan2022}. Estos antecedentes respaldan el uso de variables disponibles antes del cierre, aunque el objetivo del presente trabajo es el monto agregado de contrataciones y no la clasificación individual de oportunidades.

\subsection{Agregación temporal}

La agregación temporal transforma una serie de alta frecuencia en observaciones correspondientes a períodos más extensos. Esta operación atenúa variaciones de alta frecuencia y puede facilitar la identificación de tendencias y ciclos. La combinación de múltiples niveles de agregación permite mejorar la estimación de componentes estructurales de una serie \citep{Kourentzes2014}. También se observó que la agregación puede mejorar la precisión y reducir la incertidumbre en problemas de demanda, aunque disminuye la cantidad de observaciones disponibles \citep{Kourentzes2017}.

Este compromiso es especialmente relevante para el caso estudiado. La frecuencia semanal proporciona más puntos de entrenamiento, pero puede exhibir una variabilidad elevada. La frecuencia mensual reduce parte de ese ruido, aunque deja una muestra pequeña, mientras que el horizonte trimestral coincide mejor con la decisión financiera. Por lo tanto, la granularidad constituye una dimensión que debe evaluarse empíricamente.

El enfoque adoptado compara referencias simples, el pronóstico del CRM, modelos estadísticos, algoritmos basados en árboles y una red LSTM bajo un protocolo temporal común. La contribución no reside en proponer un algoritmo nuevo, sino en construir un proceso reproducible que transforma estados históricos del \textit{pipeline} en una proyección agregada, evita el uso de información futura y permite comparar familias de modelos sobre las mismas fechas y métricas.

\section{Objetivos}
\label{sec:objetivos}

El objetivo general del trabajo es desarrollar y evaluar un modelo de aprendizaje automático, integrado en un flujo de procesamiento de datos reproducible, para proyectar las contrataciones del trimestre siguiente a partir del estado mensual histórico de un CRM, con corrección temporal y un desempeño medido frente al pronóstico utilizado por la organización.

Dentro de los objetivos específicos se pueden mencionar los siguientes:

\begin{itemize}
    \item Diseñar e implementar un flujo reproducible para ingerir, validar, limpiar y consolidar los \textit{snapshots} históricos del CRM.
    \item Definir la variable objetivo de contrataciones y construir las variables predictoras con criterio \textit{point in time}, de modo de evitar duplicaciones y fuga de información.
    \item Analizar el efecto de las granularidades semanal, mensual y trimestral sobre la variabilidad de la serie y la precisión de los pronósticos.
    \item Comparar modelos de referencia, ARIMA, Prophet, \textit{Random Forest}, XGBoost y LSTM mediante una metodología común de validación temporal y métricas homogéneas.
    \item Seleccionar y ajustar el modelo de mejor desempeño y cuantificar su diferencia respecto del pronóstico del CRM.
    \item Generar proyecciones trimestrales agregadas, trazables e interpretables que puedan utilizarse como apoyo en la planificación financiera.
\end{itemize}

\section{Alcance del trabajo}
\label{sec:alcanceTrabajo}

El trabajo comprende el desarrollo de un prototipo funcional orientado a demostrar la factibilidad técnica y el valor práctico del enfoque. El alcance incluye la carga y validación de archivos históricos exportados desde el CRM, la consolidación de los \textit{snapshots}, la construcción de una tabla mensual de modelado, la definición de las contrataciones, la generación de variables, el entrenamiento y comparación de diferentes familias de modelos, el registro de experimentos y la exportación de proyecciones y métricas.

Quedaron fuera del alcance la conexión en tiempo real con el CRM, el desarrollo de una interfaz de usuario, el despliegue productivo y el modelado financiero posterior a la proyección. Tampoco se realizó una segmentación geográfica, debido a que la fuente disponible no contiene una dimensión de región.

% Chapter 2

\chapter{Introducción específica}

\label{Chapter2}

En este capítulo se presentan los datos, los modelos, las métricas y los recursos tecnológicos utilizados en el trabajo. Primero se caracteriza la fuente histórica proveniente del sistema de gestión comercial y se describen sus variables, su calidad y sus limitaciones. Luego se define la salida esperada y se resumen los enfoques de pronóstico considerados. Finalmente, se detallan las métricas de evaluación, las bibliotecas y la infraestructura de cómputo.

\section{Conjunto de datos}
\label{sec:conjuntoDatos}

La fuente principal de información fue una exportación histórica del CRM de una empresa del portafolio de OceanSound Partners. El archivo contenía fotografías sucesivas del embudo comercial. Cada fotografía, denominada \textit{snapshot}, registraba el estado conocido de las oportunidades en una fecha determinada. Una misma oportunidad podía aparecer en numerosos períodos consecutivos, por lo que la cantidad de filas no representaba la cantidad de operaciones independientes.

La extracción reducida reunió 1 993 467 filas y 23 columnas, con un tamaño aproximado de 498 MB. El período observado se extendió desde el 8 de enero de 2024 hasta el 1 de junio de 2026 y comprendió 126 \textit{snapshots} separados por una cadencia exacta de siete días. En total se identificaron 24 828 oportunidades, 4.476 cuentas y 225 responsables comerciales. Cada oportunidad apareció, en promedio, en 80,29 fotografías.

La tabla~\ref{tab:resumenFuente} resume las principales dimensiones de la fuente.

\begin{table}[ht]
    \centering
    \caption{Resumen del conjunto de datos utilizado en el trabajo.}
    \label{tab:resumenFuente}
    \small
    \begin{tabular}{p{7.2cm} r}
        \toprule
        Propiedad & Valor \\
        \midrule
        Filas & 1 993 467 \\
        Columnas & 23 \\
        Tamaño aproximado del archivo & 498 MB \\
        Fotografías semanales & 126 \\
        Primera fecha observada & 8 de enero de 2024 \\
        Última fecha observada & 1 de junio de 2026 \\
        Oportunidades únicas & 24 828 \\
        Cuentas únicas & 4 476 \\
        Responsables comerciales únicos & 225 \\
        Apariciones promedio por oportunidad & 80,29 \\
        \bottomrule
    \end{tabular}
\end{table}

\subsection{Estructura y variables disponibles}

La fuente tuvo una estructura de panel longitudinal. La fecha de cada fotografía indicó el instante en el que la información se encontraba disponible y el identificador de oportunidad permitió seguir una misma operación a lo largo del tiempo. Las columnas restantes describieron el estado comercial, las fechas relevantes, los montos, la probabilidad de cierre, la categoría de pronóstico y los responsables asociados.

Las variables monetarias incluyeron el valor contractual anual, los servicios profesionales y los componentes de equipamiento o licencias. También se dispuso del monto total de la oportunidad, de componentes esperados y del ingreso esperado registrado por el CRM. Los importes ya se encontraban expresados en dólares estadounidenses, por lo que no fue necesario realizar conversiones monetarias.

El estado de una oportunidad se representó mediante indicadores de cierre y resultado. Estos campos se consideraron más confiables que el nombre textual de la etapa, porque este último presentaba variantes semánticas e inconsistencias. La fecha prevista de cierre, la fecha de creación, la procedencia del contacto, el tipo de oportunidad y la categoría del pronóstico completaron la descripción comercial disponible.

La fuente contenía información suficiente para observar dos aspectos diferentes. Por una parte, permitía reconstruir el volumen y la valoración de las oportunidades que se encontraban abiertas en una fecha. Por otra, permitía identificar las operaciones que finalmente habían sido ganadas y asignarlas a su fecha efectiva de cierre. La forma en que ambas perspectivas se consolidaron sin duplicaciones ni uso de información futura se explica en el capítulo~\ref{Chapter3}.

\subsection{Contrataciones históricas}

El monto contratado se definió como la suma del valor contractual anual, los servicios profesionales y los componentes de equipamiento o licencias de las oportunidades ganadas. Se comprobó que esta suma coincidía con el monto total registrado, con una tolerancia de un dólar, para la totalidad de las operaciones consolidadas.

La historia disponible comprendió 11.983 oportunidades ganadas con fechas de cierre entre el 3 de enero de 2023 y el 29 de mayo de 2026. El monto acumulado fue de USD 186.782.760. Los servicios profesionales representaron el 73,3 por ciento, el valor contractual anual el 17,7 por ciento y los componentes de equipamiento o licencias el 9,0 por ciento restante. Esta composición mostró que la mayor parte del volumen provenía de servicios, aunque se preservó el desglose para facilitar su análisis posterior.

La cantidad de operaciones y sus montos variaron de manera considerable entre períodos. Un conjunto reducido de oportunidades de valor elevado podía modificar el total de un mes. Por este motivo, el análisis debía considerar tanto el error monetario como el error relativo y no podía apoyarse únicamente en el porcentaje de aciertos individuales.

\subsection{Calidad, limitaciones y confidencialidad}

La calidad general de las variables centrales resultó adecuada. Los valores faltantes se concentraron principalmente en la procedencia del contacto, con un 2,36 por ciento. Los identificadores de responsable y el tipo de oportunidad presentaron proporciones inferiores al 0,15 por ciento. También se detectaron 66 filas sin contenido útil en los campos principales.

La principal limitación fue la longitud temporal efectiva. Los casi dos millones de registros correspondían a repeticiones de oportunidades observadas durante menos de dos años y medio. En consecuencia, la cantidad de ejemplos útiles para comparar pronósticos temporales fue mucho menor que el número de filas originales. Esta característica condicionó la complejidad admisible de los modelos y la solidez de las estimaciones sobre períodos recientes.

Los últimos meses presentaron además censura por la derecha. Algunas oportunidades continuaban abiertas al momento de realizar la extracción y todavía podían cambiar de estado. La fuente tampoco incluyó una variable de región, por lo que no permitió evaluar pronósticos geográficos sin incorporar información externa.

Los datos crudos se procesaron dentro de la infraestructura autorizada por la organización y no se incorporaron al repositorio académico. La memoria utiliza únicamente conteos, proporciones, métricas y resultados agregados. No se publican nombres de personas, clientes u oportunidades, ni valores que permitan reconstruir una operación individual.

\section{Salida esperada e información predictiva}
\label{sec:salidaInformacion}

La salida esperada es el monto total de contrataciones del trimestre siguiente a cada origen mensual. Si $B_{T+k}$ representa el monto contratado en el mes situado $k$ períodos después del origen $T$, el objetivo trimestral $q_T$ se expresa mediante la ecuación~\ref{eq:objetivoTrimestral}.

\begin{equation}
    \label{eq:objetivoTrimestral}
    q_T = \sum_{k=1}^{3} B_{T+k}.
\end{equation}

Esta definición coincide con el horizonte utilizado por la planificación financiera y permite comparar la estimación automática con el pronóstico comercial. La observación de entrada representa la información disponible en el origen, mientras que la salida corresponde a los tres meses posteriores.

La información predictiva se organizó en tres grupos conceptuales. El primero describió el embudo comercial abierto mediante la cantidad de oportunidades, su monto total, su valor ponderado por probabilidad y el ingreso esperado. El segundo representó la historia de las contrataciones mediante valores anteriores y resúmenes móviles. El tercero describió la posición temporal de cada observación dentro del año.

\section{Modelos considerados}
\label{sec:modelosConsiderados}

El trabajo consideró modelos con diferentes supuestos y niveles de complejidad. Se incluyeron referencias simples para determinar si los métodos más elaborados aportaban una mejora real. Una referencia repitió el último valor observado y otra utilizó el valor correspondiente al mismo período del año anterior. El pronóstico ponderado del CRM se mantuvo como referencia de negocio.

Entre los modelos de series temporales se utilizaron ARIMA y Prophet. ARIMA representó la dependencia histórica mediante componentes autorregresivos, de integración y de media móvil \citep{BoxJenkins2015}. Prophet aportó una formulación descomponible de tendencia y estacionalidad \citep{TaylorLetham2018}. Ambos métodos utilizaron la historia agregada de contrataciones como entrada.

Los modelos tabulares fueron \textit{Random Forest} y XGBoost. El primero combinó múltiples árboles de regresión construidos con aleatoriedad en observaciones y variables \citep{Breiman2001}. El segundo utilizó árboles incorporados secuencialmente mediante \textit{boosting} y mecanismos de regularización \citep{ChenGuestrin2016}. Estos enfoques permitieron combinar la historia de contrataciones con variables agregadas del embudo comercial.

También se incluyó una red recurrente LSTM (\textit{Long Short Term Memory}) como alternativa para representar dependencias temporales \citep{HochreiterSchmidhuber1997}. La diversidad de métodos permitió contrastar referencias de baja complejidad, modelos estadísticos, algoritmos de árboles y una red neuronal.

La tabla~\ref{tab:modelosBibliotecas} resume las familias consideradas y sus implementaciones.

\begin{table}[ht]
    \centering
    \caption{Familias de modelos y bibliotecas utilizadas.}
    \label{tab:modelosBibliotecas}
    \small
    \begin{tabular}{p{3.4cm} p{3.9cm} p{4.9cm}}
        \toprule
        Familia & Implementación & Información principal \\
        \midrule
        Referencias simples & pandas y NumPy & Historia de contrataciones \\
        Pronóstico del CRM & Estimador propio & Embudo comercial ponderado \\
        ARIMA & statsmodels & Historia de contrataciones \\
        Prophet & Prophet & Historia de contrataciones \\
        Random Forest & scikit learn & Historia y embudo comercial \\
        XGBoost & XGBoost & Historia y embudo comercial \\
        LSTM & PyTorch & Secuencia de contrataciones \\
        \bottomrule
    \end{tabular}
\end{table}

\section{Métricas de evaluación}
\label{sec:metricasEvaluacion}

La evaluación utilizó métricas absolutas, cuadráticas, porcentuales y de ajuste. Para $n$ observaciones, donde $y_t$ es el monto real y $\hat{y}_t$ es el monto proyectado, el error absoluto medio o MAE (\textit{Mean Absolute Error}) y la raíz del error cuadrático medio o RMSE (\textit{Root Mean Squared Error}) se definen en las ecuaciones~\ref{eq:mae} y~\ref{eq:rmse}.

\begin{equation}
    \label{eq:mae}
    \mathrm{MAE} = \frac{1}{n}\sum_{t=1}^{n}\left|y_t-\hat{y}_t\right|.
\end{equation}

\begin{equation}
    \label{eq:rmse}
    \mathrm{RMSE} = \sqrt{\frac{1}{n}\sum_{t=1}^{n}\left(y_t-\hat{y}_t\right)^2}.
\end{equation}

El MAE expresa el desvío promedio en unidades monetarias y mantiene una interpretación directa. El RMSE asigna mayor peso a los errores grandes debido al término cuadrático. Esta propiedad resulta útil cuando se desea identificar modelos sensibles a meses con desvíos extraordinarios.

El error porcentual absoluto medio o MAPE (\textit{Mean Absolute Percentage Error}) expresa el error medio como porcentaje de cada observación, pero no se encuentra definido cuando el valor real es cero y puede crecer de manera desproporcionada ante valores pequeños \citep{HyndmanKoehler2006}. El error porcentual absoluto ponderado o WAPE (\textit{Weighted Absolute Percentage Error}) divide la suma de los errores absolutos por el volumen real acumulado. Ambas métricas se definen en las ecuaciones~\ref{eq:mape} y~\ref{eq:wapeChapter2}.

\begin{equation}
    \label{eq:mape}
    \mathrm{MAPE} = \frac{100}{n}\sum_{t=1}^{n}\left|\frac{y_t-\hat{y}_t}{y_t}\right|.
\end{equation}

\begin{equation}
    \label{eq:wapeChapter2}
    \mathrm{WAPE} = 100\,\frac{\sum_{t=1}^{n}\left|y_t-\hat{y}_t\right|}
    {\sum_{t=1}^{n}\left|y_t\right|}.
\end{equation}

El WAPE evita la división individual por valores nulos y pondera cada período de acuerdo con su aporte al total. Por esta razón fue la métrica principal del trabajo.

El coeficiente de determinación $R^2$ complementó las métricas de error y se calculó mediante la ecuación~\ref{eq:coeficienteDeterminacion}.

\begin{equation}
    \label{eq:coeficienteDeterminacion}
    R^2 = 1 - \frac{\sum_{t=1}^{n}\left(y_t-\hat{y}_t\right)^2}
    {\sum_{t=1}^{n}\left(y_t-\bar{y}\right)^2},
\end{equation}

donde $\bar{y}$ representa el promedio de los valores reales. Un valor cercano a uno indica que el modelo representa una proporción elevada de la variabilidad observada. Un valor negativo señala que su desempeño es inferior al de una predicción constante basada en el promedio.

\section{Entorno tecnológico}
\label{sec:entornoTecnologico}

El trabajo se desarrolló en Python 3.11. Pandas y NumPy se utilizaron para la lectura, la transformación tabular, las agregaciones temporales y los cálculos numéricos. PyArrow permitió almacenar las tablas procesadas en un formato columnar y Pandera se empleó para declarar controles sobre el esquema de entrada.

scikit learn proporcionó los modelos y contratos comunes de aprendizaje automático. XGBoost se utilizó para el modelo de \textit{boosting}, statsmodels para ARIMA, Prophet para el modelo descomponible y PyTorch para la red LSTM. MLflow permitió registrar los parámetros y las métricas de los experimentos en una base SQLite local.

Matplotlib y Seaborn se utilizaron para producir gráficos. Joblib permitió serializar el modelo seleccionado y un archivo JSON conservó sus metadatos. Pytest se empleó para las pruebas automatizadas y Ruff para el control estático del código. La tabla~\ref{tab:entornoTecnologico} resume el propósito de las herramientas principales.

\begin{table}[ht]
    \centering
    \caption{Herramientas principales del entorno tecnológico.}
    \label{tab:entornoTecnologico}
    \small
    \begin{tabular}{p{4cm} p{8.5cm}}
        \toprule
        Herramienta & Uso dentro del trabajo \\
        \midrule
        Python, pandas y NumPy & Procesamiento tabular y cálculo numérico \\
        PyArrow y Pandera & Persistencia columnar y validación del esquema \\
        scikit learn y XGBoost & Modelos tabulares \\
        statsmodels y Prophet & Modelos de series temporales \\
        PyTorch & Red neuronal recurrente \\
        MLflow y SQLite & Registro local de experimentos \\
        Matplotlib y Seaborn & Generación de gráficos \\
        Joblib y JSON & Persistencia del modelo y sus metadatos \\
        pytest y Ruff & Pruebas automatizadas y control estático \\
        \bottomrule
    \end{tabular}
\end{table}

La organización interna del código, la secuencia de procesamiento y los mecanismos de persistencia forman parte de la implementación. Por ese motivo, se desarrollan en el capítulo~\ref{Chapter3} y no se detallan en esta introducción específica.

\section{Infraestructura de cómputo}
\label{sec:infraestructuraComputo}

El desarrollo y los experimentos se ejecutaron de forma local sobre un equipo con arquitectura Apple Silicon y 24 GiB de memoria. Los modelos basados en árboles y los métodos estadísticos utilizaron el procesador central. La implementación de LSTM pudo aprovechar la aceleración provista por \textit{Metal Performance Shaders} cuando ese recurso se encontraba disponible.

Después de optimizar los tipos de datos, la tabla cargada ocupó aproximadamente 228 MB de memoria. El volumen de la extracción pudo procesarse sin infraestructura distribuida ni servicios externos de cómputo. Este entorno fue suficiente para el desarrollo y la evaluación del prototipo.

El alcance no incluyó un servicio de inferencia, una interfaz de usuario ni una conexión directa con el CRM. Las salidas del trabajo consistieron en archivos de proyección, métricas, figuras y un modelo serializado. En consecuencia, la infraestructura descrita corresponde a un entorno de desarrollo y evaluación. 
% Chapter 3

\chapter{Diseño e implementación}

\label{Chapter3}

En este capítulo se describe el diseño adoptado para transformar las fotografías históricas del CRM en una proyección del monto de contrataciones del trimestre siguiente. Primero se presenta la arquitectura del sistema y la secuencia general de procesamiento. Luego se explican la carga y validación de los datos, la consolidación temporal, la construcción de variables y la preparación del problema supervisado. Finalmente, se detallan la implementación de los modelos, el procedimiento de validación, el registro de experimentos, la generación de artefactos y las pruebas utilizadas para verificar la corrección del sistema. Los resultados cuantitativos de estas decisiones se reservan para el capítulo~\ref{Chapter4}.

\section{Arquitectura del sistema}
\label{sec:arquitecturaSistema}

La solución se implementó como un paquete de Python organizado por responsabilidades. Esta separación permitió reutilizar la misma lógica desde los cuadernos de análisis, los programas ejecutables y las pruebas automatizadas. También evitó que la preparación de datos, el entrenamiento y la generación de informes dependieran de implementaciones duplicadas.

La configuración central se concentró en un único módulo. Allí se declararon las rutas, los nombres de las columnas, los componentes monetarios, las variables que podían producir fuga de información y los parámetros temporales predeterminados. De esta manera, una modificación del esquema o del horizonte no requirió cambios dispersos en diferentes partes del código.

La tabla~\ref{tab:arquitecturaModulos} resume las capas principales de la implementación.

\begin{table}[ht]
    \centering
    \caption{Organización funcional del sistema implementado.}
    \label{tab:arquitecturaModulos}
    \small
    \begin{tabular}{p{3.2cm} p{3.4cm} p{5.1cm}}
        \toprule
        Responsabilidad & Módulos principales & Función \\
        \midrule
        Configuración & \texttt{config.py} & Centralización de rutas, columnas y reglas temporales \\
        Ingesta & \texttt{ingestion} & Lectura, tipado, validación y consolidación de fotografías \\
        Construcción de variables & \texttt{features} & Agregación mensual, rezagos, ventanas móviles y objetivo trimestral \\
        Modelos & \texttt{models} & Registro de referencias, modelos estadísticos, árboles y red recurrente \\
        Validación y entrenamiento & \texttt{validation} y \texttt{training} & Particiones temporales, métricas, comparación y ajuste \\
        Proyección & \texttt{projection} & Entrenamiento final y estimación del trimestre siguiente \\
        Informes y persistencia & \texttt{reporting}, \texttt{tracking} y \texttt{serving} & Exportación, trazabilidad y almacenamiento del modelo \\
        \bottomrule
    \end{tabular}
\end{table}

El flujo completo se expuso mediante seis programas ejecutables. El primero construyó la tabla de modelado. El segundo ejecutó la comparación mensual. El tercero realizó la comparación y la proyección trimestral. El cuarto ajustó hiperparámetros. El quinto produjo figuras y archivos para herramientas de inteligencia de negocios. El sexto entrenó y serializó el modelo destinado a la proyección. Esta secuencia permitió ejecutar cada etapa por separado y conservar resultados intermedios verificables.

\section{Ingesta y validación de los datos}
\label{sec:ingestaValidacion}

La entrada consistió en un archivo CSV exportado desde el CRM. La función de carga asignó tipos explícitos antes de construir las tablas analíticas. Los importes se representaron con variables numéricas de precisión doble, los identificadores y categorías se cargaron como variables categóricas, las fechas se interpretaron durante la lectura y los indicadores lógicos se convirtieron a un tipo booleano que admitía valores faltantes.

Los nombres originales de las columnas incluían un sufijo agregado por el mecanismo de exportación. Ese sufijo se eliminó de manera uniforme después de la lectura. También se reconocieron como ausentes las representaciones textuales habituales de valores nulos. Por último, se descartaron las filas en las que el identificador, la etapa y los principales importes se encontraban simultáneamente vacíos. Esta regla eliminó registros sin contenido útil sin aplicar una imputación general sobre variables comerciales.

La validación del esquema se implementó con Pandera. Se exigió la presencia de los identificadores y fechas indispensables, se controló que la probabilidad estuviera comprendida entre cero y cien y se rechazaron importes negativos. La categoría del pronóstico se restringió a los valores habilitados por el proceso comercial. El esquema aceptó columnas adicionales para que una ampliación no disruptiva de la exportación no interrumpiera el procesamiento. En cambio, una ausencia o una modificación incompatible de las variables centrales provocó un error temprano y detallado.

La validación se mantuvo como una opción de la función de carga. Esta decisión permitió activarla en ejecuciones de control y evitar su costo cuando se trabajaba con una fuente ya verificada. La importación de la biblioteca también se realizó de manera diferida, de modo que el núcleo de lectura pudiera utilizarse aun cuando el componente de validación no estuviera instalado.

\section{Consolidación temporal}
\label{sec:consolidacionTemporal}

El panel original contenía una fotografía por semana y repetía una misma oportunidad mientras permanecía registrada en el CRM. Para reducir la variabilidad de alta frecuencia se conservó la primera fotografía disponible de cada mes. Esta regla produjo un origen temporal mensual estable y mantuvo una relación directa con el horizonte de planificación trimestral.

La consolidación separó dos perspectivas. La primera reconstruyó las contrataciones efectivas. La segunda describió el embudo comercial que seguía abierto en cada origen. Esta separación fue necesaria porque una oportunidad ganada contribuía a la variable objetivo una sola vez, mientras que una oportunidad abierta podía formar parte de las variables predictoras en varios orígenes consecutivos.

Para construir las contrataciones se seleccionaron las oportunidades identificadas como ganadas por los indicadores del CRM. No se utilizó el texto de la etapa porque presentaba variantes semánticas. Luego se ordenaron los registros por fecha de fotografía y se conservó una sola fila por identificador de oportunidad. El monto individual $b_i$ se calculó mediante la ecuación~\ref{eq:montoContratacionImplementacion}.

\begin{equation}
    \label{eq:montoContratacionImplementacion}
    b_i = \mathrm{ACV}_i + \mathrm{PS}_i + \mathrm{HWSW}_i.
\end{equation}

Los valores ausentes de cada componente se trataron como cero exclusivamente para esta suma. Las oportunidades sin fecha de cierre se excluyeron porque no podían asignarse a un período. Después de la deduplicación, los importes se agregaron por mes de cierre y se conservaron tanto el total como el desglose por tipo de ingreso.

El embudo comercial se calculó de manera independiente para cada fotografía. Se consideró abierta una oportunidad cuyo indicador de cierre fuera falso. Sobre ese subconjunto se generaron la cantidad de oportunidades, el monto abierto, el ingreso esperado informado por el CRM y el monto ponderado por probabilidad. Para una fotografía tomada en $T$, este último se definió según la ecuación~\ref{eq:pipelinePonderado}.

\begin{equation}
    \label{eq:pipelinePonderado}
    P_T = \sum_i A_{i,T}\frac{p_{i,T}}{100},
\end{equation}

donde la suma se extiende sobre las oportunidades abiertas en $T$, $A_{i,T}$ es el monto de la oportunidad y $p_{i,T}$ es la probabilidad registrada en esa fecha. Todas las cantidades se calcularon dentro de la fotografía correspondiente, sin reemplazar sus valores por estados observados posteriormente.

También se implementó un conjunto experimental de variables más detalladas. Estas variables representaron la antigüedad media del embudo, el monto por categoría de confianza y el monto ponderado según los meses restantes hasta la fecha esperada de cierre. Su activación se controló mediante un parámetro y no formó parte del camino predeterminado. Esta separación permitió conservar la evidencia ejecutable de los ensayos sin aumentar la complejidad de la tabla utilizada habitualmente.

\section{Construcción de la tabla de modelado}
\label{sec:construccionTablaModelado}

La tabla de modelado se indexó por fin de mes. En cada fila se combinaron las contrataciones del período, el estado del embudo disponible en el origen y las variables derivadas de la historia. La unión se realizó por índice temporal y no por identificador de oportunidad, ya que la unidad de análisis del modelo fue el período y no la operación individual.

Las variables de calendario incluyeron año, mes y trimestre. El mes también se representó mediante las transformaciones seno y coseno de la ecuación~\ref{eq:codificacionCiclicaMes}. Esta codificación preservó la cercanía entre diciembre y enero, que se perdería al representar el mes solamente con un entero.

\begin{equation}
    \label{eq:codificacionCiclicaMes}
    m_{\sin} = \sin\left(2\pi\frac{m}{12}\right),
    \qquad
    m_{\cos} = \cos\left(2\pi\frac{m}{12}\right).
\end{equation}

La historia de contrataciones se incorporó mediante rezagos de uno, dos, tres y doce períodos. Los tres primeros representaron el comportamiento reciente y el último aportó una referencia anual. Se calcularon además medias móviles de tres y doce meses y tasas de variación sobre las mismas ventanas. Los cocientes infinitos producidos por una base igual a cero se reemplazaron por valores ausentes antes de construir el conjunto supervisado.

Las columnas correspondientes a la contratación total y a su desglose no se incluyeron directamente entre las entradas. Solo se utilizaron sus valores retrasados o resumidos hasta el origen. También se excluyeron de manera explícita la fecha efectiva de cierre y los indicadores de oportunidad cerrada o ganada. Estas reglas impidieron que el modelo utilizara el resultado que debía predecir.

La tabla resultante pudo persistirse en formato Parquet o CSV. Parquet se utilizó como formato predeterminado por su representación columnar y por la conservación de tipos. La alternativa CSV se mantuvo para inspecciones manuales y para facilitar el intercambio con otras herramientas.

\section{Formulación supervisada del trimestre siguiente}
\label{sec:formulacionSupervisadaTrimestre}

El objetivo operativo fue la suma de las contrataciones de los tres meses posteriores a cada origen, definida en la ecuación~\ref{eq:objetivoTrimestral}. Para los modelos tabulares se adoptó una formulación directa. Cada fila de variables observadas en $T$ se emparejó con un único valor $q_T$ correspondiente al trimestre futuro completo. De esta manera, el modelo no necesitó predecir un mes, incorporarlo como entrada y repetir el procedimiento de forma recursiva.

La preparación descartó los orígenes sin historia suficiente para calcular todos los rezagos y ventanas. También eliminó los últimos orígenes, cuyo trimestre futuro no se encontraba completamente observado. Además, se excluyeron dos orígenes recientes afectados por censura, debido a que las oportunidades correspondientes todavía podían cambiar de estado en la fuente.

La construcción de $q_T$ se realizó mediante desplazamientos hacia el futuro aplicados únicamente a la variable de salida. Las variables de entrada conservaron el índice del origen. Esta distinción quedó encapsulada en una estructura que devolvió la tabla completa, la matriz de entrada y el vector objetivo. Así se utilizó la misma alineación durante el ajuste, la comparación y la proyección.

Para los modelos de series temporales se conservó la serie mensual de contrataciones. En cada origen se pronosticaron los tres meses siguientes y se sumaron sus valores. Por lo tanto, los modelos tabulares y los modelos de serie resolvieron el mismo objetivo trimestral mediante estrategias compatibles, aunque con mecanismos de estimación diferentes.

\section{Implementación de los modelos}
\label{sec:implementacionModelos}

Los modelos se registraron bajo dos contratos comunes. Los modelos tabulares se construyeron mediante fábricas que devolvían objetos con métodos de ajuste y predicción. Los modelos de serie se representaron mediante funciones que recibían la historia disponible y devolvían el pronóstico del horizonte solicitado. Esta interfaz permitió que el procedimiento de evaluación recorriera familias heterogéneas sin incorporar lógica específica para cada biblioteca.

Se implementaron dos referencias simples. La primera repitió el último valor observado para cada mes del horizonte. La segunda utilizó los valores correspondientes a los mismos meses del ciclo anual anterior. La referencia del CRM se modeló a partir del monto abierto ponderado por probabilidad. Sobre cada ventana de entrenamiento se estimó un factor de calibración $\alpha$ sin intercepto, como se muestra en la ecuación~\ref{eq:calibracionCRM}.

\begin{equation}
    \label{eq:calibracionCRM}
    \hat{q}_T^{\mathrm{CRM}} = \alpha P_T,
    \qquad
    \alpha = \frac{\sum_T P_T q_T}{\sum_T P_T^2}.
\end{equation}

El factor se recalculó dentro de cada ventana de entrenamiento. En consecuencia, la referencia no utilizó el valor real del período evaluado para ajustar su escala.

Los modelos tabulares se implementaron mediante \textit{Random Forest} y XGBoost. Ambos recibieron las variables de calendario, los agregados del embudo y los resúmenes históricos. Los modelos ARIMA y Prophet operaron sobre la serie mensual agregada. ARIMA utilizó una especificación autorregresiva integrada de media móvil de orden $(1,1,1)$. Prophet se configuró con estacionalidad anual y sin estacionalidades semanal ni diaria, dado que la unidad de análisis fue mensual.

La red LSTM utilizó una ventana equivalente a doce períodos, una capa recurrente y una capa lineal de salida. La serie se normalizó dentro de cada historia de entrenamiento y la predicción se transformó nuevamente a la escala monetaria. La implementación seleccionó CUDA, \textit{Metal Performance Shaders} o el procesador central según la disponibilidad. Como la red se implementó para un paso futuro, se incluyó en la comparación mensual y no en la comparación trimestral directa.

Las dependencias de mayor tamaño se importaron solamente cuando el modelo correspondiente era requerido. Si una biblioteca opcional no se encontraba disponible, el registro omitía esa alternativa y permitía ejecutar el resto del flujo. Este comportamiento facilitó las pruebas parciales y evitó que una dependencia no esencial impidiera construir los datos o evaluar otros modelos.

\section{Validación temporal y ajuste}
\label{sec:validacionTemporalAjuste}

La comparación se realizó mediante validación de ventana expansiva, denominada \textit{walk forward}. En cada origen se entrenó nuevamente el modelo con todas las observaciones anteriores y se generó una predicción para un período posterior. El origen evaluado nunca formó parte de su propio conjunto de entrenamiento. A medida que avanzó el tiempo, la ventana incorporó la información nueva sin descartar la historia previa.

Se estableció el 1 de enero de 2026 como comienzo del conjunto de evaluación fuera de tiempo, denominado \textit{holdout}. Los orígenes anteriores se utilizaron para la validación cruzada temporal y los posteriores se reservaron para la evaluación final cuando contaban con un objetivo completo. El código incluyó además una regla de embargo igual al horizonte trimestral para formalizar la última fecha admisible de entrenamiento antes de un corte.

Los modelos de serie requirieron un mínimo de doce meses de historia. Los modelos tabulares comenzaron con seis filas completas, porque los rezagos y las ventanas ya habían eliminado los orígenes sin información suficiente. Aunque los mínimos fueron diferentes, las métricas de cada comparación se calcularon sobre la intersección de fechas disponible para todos los modelos. Esta regla evitó favorecer una alternativa por haber sido evaluada sobre períodos distintos.

La métrica principal fue WAPE, definida en la ecuación~\ref{eq:wapeChapter2}. También se calcularon MAE, RMSE, MAPE y $R^2$. Cada modelo produjo predicciones y valores reales separados entre validación temporal y evaluación fuera de tiempo. Cuando no existieron suficientes observaciones completas en esta última partición, el ordenamiento se apoyó en la validación temporal sin reemplazar los valores ausentes por estimaciones artificiales.

El ajuste de hiperparámetros se realizó exclusivamente sobre los orígenes anteriores al corte. Para \textit{Random Forest} se exploraron combinaciones acotadas de profundidad, cantidad mínima de observaciones por hoja y cantidad de variables consideradas en cada división. Para XGBoost se variaron la profundidad, la tasa de aprendizaje y la cantidad de árboles. Se utilizó una búsqueda exhaustiva sobre grillas pequeñas para reducir el riesgo de adaptar en exceso la selección a una historia temporal corta. También se implementó una búsqueda aleatoria como alternativa reutilizable, aunque no formó parte del flujo principal. Todas las configuraciones se compararon mediante el mismo procedimiento expansivo y la misma métrica WAPE.

\section{Trazabilidad y generación de artefactos}
\label{sec:trazabilidadArtefactos}

Los experimentos se registraron con MLflow sobre una base SQLite local. Para cada ejecución se almacenaron el nombre y la familia del modelo, la granularidad, los parámetros temporales, las métricas y el identificador del \textit{commit} de Git. Este último permitió vincular un resultado con el estado exacto del código que lo produjo.

El registro se diseñó como una capacidad opcional. Si MLflow no se encontraba disponible o fallaba la conexión, el entrenamiento continuaba y conservaba los resultados en archivos CSV. Esta decisión evitó que una falla de observabilidad invalidara una ejecución computacionalmente costosa.

La etapa de informes generó cuatro figuras: comparación de WAPE entre modelos, evolución histórica con la proyección, valores reales frente a valores estimados y composición de las contrataciones por tipo de ingreso. Los períodos afectados por censura se diferenciaron visualmente de los períodos completos. También se produjeron tres tablas en formato largo para su utilización en herramientas de inteligencia de negocios: historia de contrataciones, tabla comparativa de modelos y proyección del trimestre siguiente.

Después de seleccionar una configuración, el entrenamiento final utilizó todas las observaciones admisibles. El estimador se almacenó con Joblib y se generó un archivo JSON con las columnas requeridas, la clase del modelo, la fecha de entrenamiento, los parámetros, la cantidad de filas y el origen de la proyección. El programa de entrenamiento final volvió a cargar el artefacto y comprobó que la predicción coincidiera con la obtenida antes de guardarlo.

\section{Pruebas de la implementación}
\label{sec:pruebasImplementacion}

La batería automatizada se concentró en los invariantes que podían alterar la validez del análisis. El repositorio incluyó cuarenta y cinco casos de prueba al considerar las parametrizaciones. Las pruebas unitarias utilizaron tablas pequeñas y resultados conocidos para detectar errores que serían difíciles de identificar sobre el conjunto completo.

La tabla~\ref{tab:pruebasImplementacion} resume las verificaciones principales.

\begin{table}[ht]
    \centering
    \caption{Cobertura funcional de las pruebas automatizadas.}
    \label{tab:pruebasImplementacion}
    \small
    \begin{tabular}{p{3.5cm} p{8.4cm}}
        \toprule
        Área & Invariante verificado \\
        \midrule
        Esquema de entrada & Rechazo de categorías inválidas, probabilidades fuera de rango e importes negativos \\
        Consolidación & Conteo único de oportunidades ganadas y selección mensual correcta \\
        Corrección temporal & Invariancia de una fotografía al agregar datos futuros y exclusión de columnas posteriores al cierre \\
        Variables & Cálculo de rezagos, ventanas, tasas y alineación estrictamente futura del objetivo \\
        Validación & Ventana expansiva sin superposición y separación correcta del corte temporal \\
        Comparación & Uso de fechas comunes para las métricas y ordenamiento consistente de alternativas \\
        Proyección & Suma correcta del trimestre futuro y tratamiento de los orígenes censurados \\
        Persistencia & Equivalencia de las predicciones antes y después de guardar el modelo \\
        \bottomrule
    \end{tabular}
\end{table}

La corrección \textit{point in time} recibió pruebas específicas. Una de ellas calculó las variables de una fotografía con el panel completo y volvió a calcularlas después de eliminar todas las observaciones futuras. Ambos resultados debían coincidir. Otras pruebas verificaron que una entrada en $T$ se asociara siempre con una salida posterior, que las columnas de resultado no aparecieran entre las variables predictoras y que el entrenamiento de cada origen terminara antes de la observación evaluada.

En síntesis, la implementación convirtió un panel semanal con oportunidades repetidas en una tabla mensual apta para pronosticar el trimestre siguiente. La separación entre ingesta, variables, modelos, validación y persistencia permitió aplicar las mismas reglas temporales en todas las alternativas. La formulación directa, la evaluación expansiva y las pruebas contra fuga de información proporcionaron una base reproducible para los ensayos que se presentan en el capítulo~\ref{Chapter4}.

\include{Chapters/Chapter4}
\include{Chapters/Chapter5}

%----------------------------------------------------------------------------------------
% Apéndices
%----------------------------------------------------------------------------------------

\appendix

% Incluir apéndices desde archivos separados si es necesario
%\include{Appendices/AppendixA}

%----------------------------------------------------------------------------------------
% Bibliografía
%----------------------------------------------------------------------------------------

\renewcommand{\bibname}{Bibliografía} % Para asegurarte de que el título sea correcto
\phantomsection % Necesario para que el enlace del marcador sea correcto

\printbibliography[heading=bibintoc]

\end{document}






